\chapter{기계의 입력: 눈, 귀, 그 밖의 센서는 어떻게 연결되는가}
\chaptermark{기계의 입력}
\label{sec:sensors}
\begin{chaptergoals}
\item 변환기로서의 감각 기관: 수용체 게이트, 이득 조절, 스파이크 인코더;
\item 광자의 산탄 잡음이 어둠 속 시각을 제한하는 이유와 어스름에 더 느리게 보는 이유;
\item 자동 이득 조절이 거대한 입력 범위를 스파이크에 담아 센서가 변화를 알리게 하는 방식;
\item 눈이 영상을 백 배로 압축하는 방식과 귀가 필터 뱅크로 일하는 방식.
\end{chaptergoals}

\section{변환기로서의 센서}

그림~\ref{fig:machine}에서 데이터는 왼쪽의 “센서” 블록으로부터 기계에 들어왔다. 이제 그 블록을 열어 보자. 엔지니어에게 모든 감각 기관은 아날로그 입력단과 끝에 아날로그-디지털 변환기를 갖춘 \emph{변환기}이다. 다만 출력되는 디지털 값은 숫자가 아니라 스파이크이다.

감각 기관의 회로는 모두 같다(그림~\ref{fig:transducer}). 빛, 압력, 진동, 분자 같은 물리량이 감각 세포 막의 수용체 단백질에 작용한다. 이 단백질은 게이트처럼 동작한다. 스스로, 또는 짧은 반응 사슬을 거쳐 이온 채널을 연다. \emph{수용체 전류}가 생기고, 그와 함께 막에 아날로그 전압이 생긴다. 이 전압은 이득 조절을 거쳐 스파이크 인코더에 들어간다. 인코더는 이미 익숙한 적분 뉴런~\eqref{eq:glutneuron}으로, 수준을 발화율과 스파이크 시각으로 바꾼다. 출력은 거의 언제나 같다. 눈, 귀, 후각, 피부의 감각 뉴런은 글루탐산을 방출하므로, 입력은 곧바로 \nt{GLUT} 버스에 연결된다.

\begin{figure}[t]
\centering
\begin{tikzpicture}[>=Stealth,
  blk/.style={draw,very thick,rounded corners=2pt,align=center,font=\footnotesize,minimum height=1.0cm,text width=1.9cm,fill=blue!5}]
\node[align=center,font=\small] (P) at (0.1,0) {빛, 소리,\\압력,\\분자};
\node[blk] (R) at (3.0,0) {수용체\\게이트};
\node[blk] (G) at (6.5,0) {이득\\조절};
\node[blk] (E) at (10.0,0) {스파이크\\인코더};
\node[align=center,font=\small] (B) at (13.4,0) {\nt{GLUT}\\버스};
\draw[->,very thick] (P) -- (R);
\foreach \ic/\xx in {sun/-1.1,speaker/-0.3,press/0.5,molecule/1.3}{\pic[scale=0.85] at (\xx,1.05) {\ic};}
\draw[->,very thick] (R) -- node[above,font=\scriptsize] {전류} (G);
\draw[->,very thick] (G) -- node[above,font=\scriptsize] {수준} (E);
\draw[->,very thick,red!70!black] (E) -- node[above,font=\scriptsize,black] {스파이크} (B);
\draw[decorate,decoration={brace,amplitude=5pt,mirror},gray!60!black] (1.8,-0.8) -- (7.7,-0.8) node[midway,below=5pt,font=\scriptsize,black] {아날로그 부분};
\draw[decorate,decoration={brace,amplitude=5pt,mirror},gray!60!black] (8.8,-0.8) -- (11.2,-0.8) node[midway,below=5pt,font=\scriptsize,black] {스파이크};
\end{tikzpicture}
\caption{변환의 사슬로서의 감각 기관. 수용체 단백질이 채널을 열어 전류를 만들고, 이득 조절이 그것을 환경에 맞추며, 인코더~\eqref{eq:glutneuron}가 수준을 스파이크로 바꿔 \nt{GLUT} 버스로 보낸다. 눈과 귀의 첫 단들은 스파이크를 전혀 내지 않는다. 신호를 멀리 보내야 할 때까지 신호는 아날로그로 남는다.}
\label{fig:transducer}
\end{figure}

전자 엔지니어에게 익숙한 세부가 하나 있다. 눈과 귀에서 맨 처음 세포들은 스파이크를 전혀 내지 않는다. 보드 위의 아날로그 단처럼 이웃에게 매끄러운 전압을 넘긴다. 스파이크는 신호를 멀리 날라야 하는 곳에서만 나타난다. 아날로그 신호는 수 밀리미터만 가도 감쇠하고 잡음이 끼지만, 스파이크는 \ref{sec:neurontypes}장에서 보았듯이 같은 높이로 전선 끝까지 간다. 디지털화는 긴 선로가 시작되는 곳에서 일어난다.

\section{물리의 한계에서}

뇌의 센서는 감도의 물리적 한계 근처에서 동작한다. 어둠 속에서 빛 센서는 \emph{광자 하나}에 응답한다. 빛 입자 하나를 흡수하면 약 1피코암페어의 전류가 생기는데, 이는 센서 자체의 잡음 위로 드러난다~\cite{Baylor1979}. 소리 센서는 자기 감지 다발이 1나노미터보다 작게 움직여도 알아챈다~\cite{Hudspeth1989}.

빛이 적을 때 정밀도를 제한하는 것은 센서가 아니라 빛 자체이다. 광자는 무작위로, 푸아송 흐름으로 도착한다. 적분 시간 $T$ 동안 센서에 평균 $N=\Phi T$개의 광자가 들어오면, 그 수는 $\sqrt N$만큼 요동한다. 증가분 $\Delta N$이 눈에 띄려면 이 요동보다 몇 배 커야 하므로, 구별할 수 있는 밝기 변화의 비율은 다음과 같다.
\begin{empheq}[box=\keybox]{equation}
\frac{\Delta I}{I} \;\approx\; \frac{k}{\sqrt{\Phi\,T}} .
\label{eq:photonnoise}
\end{empheq}
이것은 스파이크 수에 대한 식~\eqref{eq:rateerror}와 같은 식이다. 광자의 산탄 잡음과 스파이크의 산탄 잡음은 같은 법칙을 따른다. 어둠 속에서 눈이 정밀도를 높이는 방법은 광자를 더 오래 모으는 것 하나뿐이고, 실제로 적분 시간은 수십분의 1초까지 늘어난다. 그래서 어스름 속에서는 우리가 더 느리게 본다.

\section{동적 범위: 이득 조절}

센서의 가장 어려운 공학 문제는 범위이다. 조도는 별이 뜬 밤에서 햇빛 비치는 눈밭까지 약 열 자릿수 변하고, 소리의 세기는 가청 문턱에서 통증 문턱까지 열두 자릿수 변한다. 반면 스파이크 발화율은 0에서 초당 수백까지, 세 자릿수가 안 되게 변한다. 이런 입력을 이런 출력에 선형으로 담을 수는 없다.

해법은 여느 수신기와 같다. \emph{자동 이득 조절}이다. 눈의 감각 세포 응답은 다음 식으로 잘 기술된다.
\begin{empheq}[box=\keybox]{equation}
r \;=\; r_{\max}\,\frac{I}{I+\sigma} ,
\label{eq:nakarushton}
\end{empheq}
여기서 $\sigma$는 응답이 절반에 이르는 밝기이다~\cite{NakaRushton1966}. \eqref{eq:nakarushton} 자체는 두세 자릿수만 담는다. 그러나 $\sigma$는 일정하지 않다. 밝은 배경에서 오래 동작하면 평균 밝기를 따라 커진다. 응답 곡선은 밝기의 로그 축을 따라 이동하며, 매번 가파른 구간을 지금 입력이 있는 곳에 놓는다(그림~\ref{fig:adaptation}). 이것은 \ref{sec:gaba}장의 분로 억제와 같은, 전체 활동에 의한 나눗셈이다~\cite{Carandini2012}. 다만 여기서 분모는 지난 몇 초 동안의 평균 밝기이다.

\begin{figure}[t]
\centering
\begin{tikzpicture}[>=Stealth,x=1.25cm,y=2.8cm]
\draw[->,gray!70!black] (-2,0) -- (6.4,0) node[below left,font=\small,black] {$\log_{10} I$};
\draw[->,gray!70!black] (-2,0) -- (-2,1.15) node[left,font=\small,black] {$r/r_{\max}$};
\foreach \u in {-2,0,2,4,6}{\draw[gray!60!black] (\u,-0.02) -- (\u,0.02); \node[below,font=\scriptsize] at (\u,0) {$\u$};}
\draw[densely dashed,gray!60!black] (-2,0.5) -- (6.2,0.5);
\foreach \s/\c/\lab in {0/blue!60!black/{어두운 배경},2/green!45!black/{중간},4/red!70!black/{밝은 배경}}{
  \pgfmathsetmacro{\lo}{max(-2,\s-3)}
  \draw[very thick,\c] (-2,0) -- (\lo,0) plot[domain=\lo:6.2,samples=80] (\x,{1/(1+10^(\s-\x))});
  \node[font=\scriptsize,\c,anchor=west] at (\s+0.15,0.42) {\lab};}
\foreach \s/\ic in {0/moon,2/cloud,4/sun}{\pic at (\s+0.6,0.64) {\ic};}
\end{tikzpicture}
\caption{\eqref{eq:nakarushton}에 따른 빛 센서의 이득 조절. 각 곡선은 밝기의 두세 자릿수를 담는다. 더 밝은 배경으로 옮겨 가면 $\sigma$가 커지고 곡선은 로그 축을 따라 이동한다. 이동하는 곡선들이 함께 전체 범위를 덮는다.}
\label{fig:adaptation}
\end{figure}

이득 조절에는 \ref{sec:code}장에서 본 대가가 따른다. 센서는 밝기가 아니라 \emph{배경에 대한} 밝기를 알린다. 일정한 입력은 점차 응답을 일으키지 못하게 되고, 변화는 매번 응답을 일으킨다. 기계의 입력은 처음부터 수준이 아니라 변화를 말하며, 여기에 명제~\ref{prop:economy}와 같은 스파이크 절약이 있다~\cite{Barlow1961}. \ref{sec:glutcomputer}장의 켜짐 센서와 꺼짐 센서도 여기서 나온다~\cite{Schiller1992}. 한쪽은 더 밝아졌다고, 다른 쪽은 더 어두워졌다고 알린다.

엔지니어들은 이 기법을 \emph{이벤트 카메라}에서 재현했다. 이런 카메라는 고정된 주기로 프레임을 찍지 않는다. 각 픽셀이 공통 클럭 없이 스스로, 밝기의 로그가 정해진 비율만큼 변하면 “더 밝음” 또는 “더 어두움” 이벤트를 낸다~\cite{Lichtsteiner2008}. 이것은 \ref{sec:glutcomputer}장의 켜짐---꺼짐 이중 레일 부호를 실리콘으로 구현한 것 그대로이며, 일반 이미지 센서로는 얻을 수 없는 범위와 속도를 카메라에 준다~\cite{Gallego2022}.

\section{눈: 전선 하나로의 압축}

눈에는 약 $10^8$개의 빛 센서가 있지만~\cite{Curcio1990}, 영상을 뇌로 나르는 출력 뉴런은 약 백만 개이다. 신호가 눈을 떠나기 전에 약 백 배로 압축된다. 주요 압축 기법은 우리에게 익숙하다. 각 출력 뉴런은 작은 점의 밝기와 그 주변 밝기의 차이에 응답한다. \ref{sec:gaba}장에서처럼 억제를 통한 이웃 빼기이다. 영상의 고른 부분은 거의 비용이 들지 않고, 가장자리와 변화가 전송된다. 출력 뉴런은 전문화되어 있다. 어떤 것은 밝아짐을, 어떤 것은 어두워짐을, 또 어떤 것은 특정 방향의 움직임을 알린다(그림~\ref{fig:direction}). 생쥐에서는 이런 출력 채널이 서른 개 넘게 세어졌다~\cite{Baden2016}.

이렇게 만들어진 케이블의 대역폭은 얼마일까? 기니피그의 출력 뉴런 기록으로부터 $\sim10^5$개의 뉴런에서 초당 약 $875$천 비트가 측정되었다. 이를 사람의 출력 뉴런 백만 개로 환산하면 초당 $10^7$비트 정도가 된다~\cite{Koch2006}. 옛날 10메가비트 네트워크 케이블과 같다. 뉴런당 평균 발화율이 초당 스파이크 몇 개라면, 이는 스파이크당 몇 비트로, \ref{sec:code}장에서 얻은 것과 같다.

\section{귀: 필터 뱅크}

귀는 다른 문제를 푼다. 소리를 주파수별로 분해해야 한다. 엔지니어라면 대역 통과 필터 뱅크를 달 것이고, 귀는 바로 그 일을 기계적으로 한다. 소리 진동은 탄성 칸막이를 따라 달리는데, 칸막이의 성질은 길이를 따라 매끄럽게 변하고, 각 구간은 자기 주파수에 공진한다. 그 구간에 앉은 센서는 자기 대역에만 응답한다(그림~\ref{fig:filterbank}). 주파수는 위치에 따라 대략 로그로 배치된다. 옥타브마다 비슷한 몫의 센서가 할당된다. 동작한 센서의 번호가 곧 주파수이며, 이것이 \ref{sec:code}장의 위치 부호이다.

\begin{figure}[t]
\centering
\begin{tikzpicture}[>=Stealth,x=1.35cm,y=2.2cm]
\draw[->,gray!70!black] (-0.3,0) -- (7.6,0) node[above left,font=\small,black] {주파수, Hz};
\draw[->,gray!70!black] (-0.3,0) -- (-0.3,1.15) node[left,font=\small,black] {응답};
\foreach \k/\f in {0/125,1/250,2/500,3/1000,4/2000,5/4000,6/8000}{
  \draw[gray!60!black] (\k,-0.02) -- (\k,0.02); \node[below,font=\scriptsize] at (\k,0) {\f};
  \draw[thick,blue!60!black] plot[domain={max(\k-1.2,-0.3)}:{\k+0.7},samples=60] (\x,{1/(1+((\x-\k)/(\x<\k?0.45:0.2))^4)});}
% a piano keyboard: seven white keys per octave, like the octaves on the axis
\foreach \i in {0,...,48}{\draw[gray!50!black,fill=white,line width=0.3pt] ({\i/7},-0.44) rectangle ({(\i+1)/7},-0.26);}
\foreach \o in {0,...,6}{\foreach \j in {0,1,3,4,5}{\fill[black] ({\o+(\j+1)/7-0.035},-0.35) rectangle ({\o+(\j+1)/7+0.035},-0.26);}}
\end{tikzpicture}
\caption{대역 통과 필터 뱅크로서의 귀, 도식. 각 곡선은 한 구간의 센서가 여러 주파수의 소리에 보이는 응답이다. 중심 주파수는 로그 눈금 위의 건반처럼 한 옥타브 간격이다. 실제 필터는 고주파 쪽 경사가 가파르고 저주파 쪽 경사가 완만하다. 아래는 건반이다. 귀에서처럼 건반에서도 옥타브마다 같은 자리가 할당된다.}
\label{fig:filterbank}
\end{figure}

귀의 공진기는 수동적이지 않다. 센서 일부는 소리에 맞춰 스스로 칸막이를 흔들어 약한 진동을 진폭으로 수천 배($66$--$76$\,dB) 증폭하고, 소리가 커지면 증폭은 줄어든다~\cite{Ruggero1997,Robles2001}. 엔지니어에게 이것은 자기 발진 경계 바로 옆에 둔 양의 되먹임 증폭기이다. \ref{sec:gaba}장의 네트워크와 같은 경계 위의 삶이다. 경계 근처에서 시스템은 작은 신호에 특히 민감하다. 음량 압축도 같은 증폭기가 한다. 작은 소리는 크게, 큰 소리는 작게 증폭한다.

귀에는 두 번째 부호도 있다. 낮은 주파수에서 뉴런의 스파이크는 소리와 위상을 맞춰 나간다. 뉴런은 모든 파동마다는 아니지만 각 파동의 특정 구간에서 발화한다. 기니피그에서 위상 고정은 약 $600$\,Hz에서 약해지기 시작해 약 $3.5$\,kHz까지 여전히 드러난다~\cite{PalmerRussell1986}. 그래서 스파이크 시각에는 소리의 정확한 시간이 남고, 거기서 두 귀에 소리가 도착하는 시간 차가 나온다. \ref{sec:glutcomputer}장의 네트워크는 이 시간 차로 방향을 찾는다~\cite{Grothe2010}. 높은 주파수에서는 스파이크가 파동을 따라가지 못하고 위치 부호만 남는다. 귀 하나가 \ref{sec:code}장의 두 부호를 모두 쓴다. 뉴런이 따라갈 수 있는 곳에서는 시간을, 그렇지 못한 곳에서는 위치를 쓴다.

\section{나머지 센서}

\begin{table}[t]
\centering
\footnotesize
\setlength{\tabcolsep}{3pt}
\renewcommand{\arraystretch}{1.15}
\begin{tabular}{>{\raggedright}p{1.9cm}>{\raggedright}p{2.6cm}>{\raggedright}p{4.0cm}>{\raggedright\arraybackslash}p{4.6cm}}
\toprule
감각 & 무엇을 재는가 & 입력의 구조 & 이 책의 기법 \\
\midrule
시각 & 빛 & 센서 $\sim10^8$개, $\sim100$배 압축, $\sim10^7$ bit/s & 이득 조절, 이웃 빼기, 두 전선, 움직임 방향 \\
청각 & 공기 압력 & 기계적 필터 뱅크, 능동 증폭기 & 위치 부호와 타이밍 부호, 지연 \\
촉각 & 압력, 진동 & 빨리 적응하는 센서와 느리게 적응하는 센서 & “고역 통과 필터 --- 저역 통과 필터” 쌍 \\
온도 & 열 & 따뜻함과 차가움의 별도 센서 & 이중 레일 부호 \\
후각 & 분자 & 수용체 유형 $\sim400$개, 센서마다 한 유형 & 조합 부호: 냄새는 동작한 유형의 집합 \\
미각 & 음식 속 물질 & 다섯 가지 맛: 단맛, 쓴맛, 짠맛, 신맛, 감칠맛 & 전용 선로; 일부 세포는 글루탐산이 아니라 ATP나 세로토닌을 방출 \\
몸의 위치 & 근육의 신장 & 신장의 길이와 속도 센서 & 운동 조절기를 위한 되먹임 \\
\bottomrule
\end{tabular}
\caption{센서는 어떻게 연결되어 있는가. 셋째 열의 구조는 모두 측정되었다. 오른쪽 열은 이를 앞 장들의 기법과 연결한다.}
\label{tab:sensors}
\end{table}

나머지 감각은 표~\ref{tab:sensors}에 모았다. 거의 모든 행에서 앞 장들의 기법을 알아볼 수 있다. 촉각은 필터 쌍을 쓴다. 어떤 압력 센서는 닿는 순간 응답하고 곧 조용해지며, 다른 센서는 압력이 있는 동안 응답을 유지한다. 온도는 따뜻함과 차가움을 따로 맡는 두 전선이 전한다. 가장 특이한 입력은 후각이다. 사람의 후각 수용체 유형은 약 사백 가지이고, 후각 뉴런 하나는 한 유형만 지닌다~\cite{Mombaerts2004}. 냄새는 한 유형이 아니라 여러 유형의 집합을 활성화하며, 메시지는 \emph{어떤} 유형이 동작했는가이다. 엔지니어에게 이것은 길이 약 사백의 이진 벡터이고, 가능한 값의 수는 천문학적이다.

\begin{keyblock}
\begin{proposition}[기계의 입력]
\label{prop:sensors}
모든 감각은 감도의 물리적 한계 근처에서 동작하는 변환기이다. 거대한 동적 범위를 자동 이득 조절로 압축하고, \nt{GLUT} 버스로 무엇보다 \emph{변화}를 보낸다. 이를 위해 센서는 기계 자신과 같은 기법을 쓴다. 이중 레일 부호, 위치 부호, 타이밍 부호, 배경에 의한 나눗셈이다. 센서의 구조는 측정되었다. 그 부호가 스파이크당 최대 정보에 맞춰져 있다는 것은 근거가 탄탄한 가설이다.
\end{proposition}
\end{keyblock}

입력도 다른 모든 것처럼 비동기로 동작한다. 각 센서는 알릴 것이 있을 때 스파이크를 내고, 감각마다 도착 지연이 다르다. 소리는 수 밀리초, 빛은 수십 밀리초 뒤에 온다. 공통 클럭이 없는데 기계가 이것들을 어떻게 하나의 세계상으로 합치는가는 \ref{sec:synthesis}장의 시간 맞추기 문제 가운데 하나이다.

\section{연습문제}

\begin{exercise}[\lvlA]
\label{ex:11:1}
\emph{광자를 얼마나 모아야 하는가.} 어스름 속에서 센서에 초당 광자 $100$개가 들어온다. 증가분은 요동~\eqref{eq:photonnoise}의 $3$배보다 크면 눈에 띈다. (a)~센서 하나가 밝기의 $10\,\%$ 변화를 알아채는 데 얼마나 걸리는가? (b)~$0.2\,\text{s}$ 안에 해내려면 센서 몇 개를 합해야 하며, 축 방향 해상도는 몇 배 떨어지는가?
\end{exercise}

\begin{exercise}[\lvlA]
\label{ex:11:2}
\emph{스파이크 하나에 몇 비트.} (a)~눈은 초당 스파이크 $3$--$5$개로 발화하는 출력 뉴런 $10^6$개로 $\sim10^7$ bit/s를 보낸다. $\Delta t=1\ms$일 때 용량~\eqref{eq:spikeentropy}의 몇 분의 몇을 쓰는가? (b)~냄새는 $\approx400$개 수용체 유형 중 $20$개를 켠다. 이 집합은 몇 비트를 나르며, 무작위 냄새 두 개는 평균 몇 개의 유형을 공유하는가?
\end{exercise}

\begin{exercise}[\lvlB]
\label{ex:11:3}
\emph{곡선~\eqref{eq:nakarushton} 하나로 할 수 있는 것.} (a)~응답이 $r_{\max}$의 $10\,\%$에서 $90\,\%$로 커지는 밝기 범위는? 몇 자릿수인가? (b)~밝기 열 자릿수를 덮으려면 이동하는 곡선의 위치가 몇 개 필요한가? 음량 열두 자릿수라면?
\end{exercise}

\begin{exercise}[\lvlB]
\label{ex:11:4}
\emph{귀에서 타이밍 부호의 한계.} 두 귀의 도착 시간 차는 최대 $0.22\,\text{m}/343\,\text{m/s}$이다. (a)~스파이크가 순음과 위상을 맞춰 나갈 때, 어느 주파수까지 방향이 한 가지로 정해지는가? (b)~스파이크는 $0.5\ms$만큼 흔들리는데 $20\,\mu\text{s}$를 구별해야 한다. 초당 스파이크 $100$개인 뉴런 몇 개를 $50\ms$ 동안 평균해야 하는가?
\end{exercise}

\begin{exercise}[\lvlB]
\label{ex:11:5}
\emph{눈의 케이블에 이벤트 카메라를.} 출력이 $10^7$ bit/s인 $10^6$픽셀 카메라는 픽셀의 $\ln I$가 $\ln1.2$만큼 변하면 이벤트(주소와 부호)를 보낸다. 밝기 차가 $4$배이고 길이가 $500$픽셀인 가장자리는 최대 얼마나 빨리 움직일 수 있는가? 픽셀당 $8$비트인 일반 카메라는 같은 출력으로 초당 몇 프레임을 보낼 수 있는가?
\end{exercise}

\begin{exercise}[\lvlB]
\label{ex:11:6}
\emph{시뮬레이션: 이득 조절.} 센서~\eqref{eq:nakarushton}는 $\tau=1\,\text{s}$로 $\sigma$를 로그로, $\tau\,\dd\ln\sigma/\dd t=\ln I-\ln\sigma$, 또는 선형으로, $\tau\,\dd\sigma/\dd t=I-\sigma$ 조정한다. 배경은 $1\to100\to1$로 도약한다. 각 법칙에서 위로, 아래로 도약한 뒤 $+10\,\%$ 시험 자극에 대한 응답이 적응된 값의 절반까지 회복되는 데 몇 초가 걸리는가?
\end{exercise}

\begin{exercise}[\lvlC]
\label{ex:11:7}
\emph{곡선은 어떤 세계에 맞춰져 있는가.} 출력 잡음이 작고 일정할 때, 밝기에 대한 정보는 $r(I)=r_{\max}\Phi_I(I)$일 때 최대이다. $\Phi_I$는 밝기의 분포 함수이다. (a)~곡선~\eqref{eq:nakarushton}이 최적인 분포는 무엇이며, 그 $\log_{10}I$의 퍼짐은? (b)~출력 잡음이 푸아송일 때 최적인 곡선은?
\end{exercise}
